% Generated draft: requires mathematical and source review.
\documentclass[11pt]{article}
\usepackage[margin=25mm]{geometry}
\usepackage{fontspec}
\setmainfont{Latin Modern Roman}
\usepackage{amsmath,amssymb,mathtools,unicode-math}
\setmathfont{Latin Modern Math}
\usepackage{enumitem,needspace,xurl,hyperref}
\setlength{\parindent}{0pt}
\setlength{\parskip}{5pt}
\setlength{\emergencystretch}{3em}
\newcommand{\N}{\mathbb{N}}
\newcommand{\Z}{\mathbb{Z}}
\newcommand{\Q}{\mathbb{Q}}
\newcommand{\R}{\mathbb{R}}
\newcommand{\C}{\mathbb{C}}
\newcommand{\F}{\mathbb{F}}
\newcommand{\A}{\mathbb{A}}
\newcommand{\PP}{\mathbb{P}}
\newcommand{\E}{\mathbb{E}}
\newcommand{\eps}{\varepsilon}
\newcommand{\dd}{\,\mathrm{d}}
\begin{document}
\section*{UnsolvedMath ID: 800047}
\textbf{Source problem code:} catalogue entry
\subsection{Definitions and mathematical statement}
A semigroup is a nonempty set equipped with an associative binary operation. A semigroup \(S\) is periodic if every element generates a finite subsemigroup; equivalently,
\[
\forall s\in S\ \exists m,p\in\N,\quad m\geq 1,\ p\geq 1,\quad s^{m+p}=s^m.
\]
Let \(\mathcal P\) denote the class of periodic semigroups.

The supplied source uses the technical terms \emph{inclusive variety}, \emph{elementary inclusive variety}, and \emph{nonelementary inclusive variety}, but supplies no definitions of them or of the relevant lattice operations. These terms therefore remain unresolved; in particular, no identification with ordinary semigroup varieties or first-order elementary classes is assumed.

For a class \(\mathcal C\) of semigroups, write \(L(\mathcal C)\) provisionally for the source's lattice of inclusive varieties within \(\mathcal C\), whenever that lattice is defined. Write \(E(\mathcal C)\) and \(N(\mathcal C)\) for its elementary and nonelementary members, respectively, using the source's as-yet-unspecified meanings. Denote its meet and join by \(\wedge\) and \(\vee\).

The problem, attributed in the source statement to G. Mashevitzky, has two parts.
\begin{enumerate}
\item Describe the classes \(\mathcal C\) of semigroups for which \(E(\mathcal C)\) forms a sublattice of \(L(\mathcal C)\). The binary closure requirement is
\[
\forall U,V\in E(\mathcal C),\qquad
U\wedge V\in E(\mathcal C)
\quad\text{and}\quad
U\vee V\in E(\mathcal C).
\]
The source does not specify which classes \(\mathcal C\) are admissible or how the relative lattice is constructed.
\item Within \(\mathcal P\), describe the nonelementary inclusive varieties that form an upper subsemilattice of \(L(\mathcal P)\). Under the usual interpretation of an upper subsemilattice as a join-closed family, a family \(\mathcal U\subseteq N(\mathcal P)\) must satisfy
\[
\forall U,V\in\mathcal U,\qquad U\vee V\in\mathcal U.
\]
The supplied wording does not determine whether the intended target is the whole family \(N(\mathcal P)\), particular subfamilies \(\mathcal U\), or some distinguished family requiring additional hypotheses. The displayed condition records the closure requirement without choosing among these targets.
\end{enumerate}

As background, the supplied source reports that elementary inclusive varieties form a sublattice, and nonelementary inclusive varieties form an upper subsemilattice, in the lattice of inclusive varieties of finite Abelian groups, citing \([\mathrm{M3}]\). This is retained as a source-reported result; its definitions, hypotheses, and proof are not provided in the record.
\subsection{Short English statement}
G. Mashevitzky asks which classes of semigroups have elementary inclusive varieties closed under lattice meet and join, and asks for a description of nonelementary inclusive varieties forming an upper subsemilattice among periodic semigroups. The technical definitions and the precise target of the second request are missing. The source reports the corresponding closure properties for finite Abelian groups, citing \([\mathrm{M3}]\).
\subsection{Sources}
\begin{itemize}
\item Supplied local UnsolvedMath registry record: permanent ID \(800047\), source code \texttt{catalogue entry}, title \emph{The Lyapin's Notebook --- Problem 49}; upstream revision \texttt{d371ee917ac9647fbb06cd6fc529deec6f3dda82}. The registry labels the problem \emph{open}; this is supplied status metadata, not an independently verified assessment.
\item Bershadsky, Kublanovsky, and Mashevitzky, \emph{The Lyapin's Notebook: Unsolved Problems in Semigroup Theory} (2026), Problem 49. Supplied source URL: \url{https://arxiv.org/abs/2604.04763}. The background records access on 2026-10-16 and extraction by \texttt{arxiv\_latex}; these are supplied provenance details, not a claim of independent access or verification.
\item The problem statement explicitly attributes the problem to G. Mashevitzky. The registry's \texttt{proposed\_by} field lists Bershadsky, Kublanovsky, and Mashevitzky, with proposed year 2026; these attribution roles are preserved separately.
\item Reference \([\mathrm{M3}]\) is cited in the supplied statement for the finite Abelian group background result. No bibliographic entry or link identifying \([\mathrm{M3}]\) is supplied.
\end{itemize}
\end{document}
